\documentclass[11pt]{article}

% ===================================================================
% Shared preamble for the "tiny lemma, read slowly" preprint series.
% Mirrors the style of FrobeniusLemmas_preprint.tex.
% ===================================================================

\usepackage{amsmath,amssymb,amsthm}
\usepackage{mathtools}
\usepackage[margin=1.2in]{geometry}
\usepackage{hyperref}
\usepackage{xcolor}
\usepackage{listings}
\usepackage{tikz}
\usepackage{tikz-cd}
\usepackage{mathpartir}
\usepackage{newunicodechar}
\usepackage{setspace}

\onehalfspacing

% ---- Unicode glyphs ----
\newunicodechar{→}{\ensuremath{\to}}
\newunicodechar{↦}{\ensuremath{\mapsto}}
\newunicodechar{↔}{\ensuremath{\leftrightarrow}}
\newunicodechar{⟹}{\ensuremath{\Longrightarrow}}
\newunicodechar{⟶}{\ensuremath{\longrightarrow}}
\newunicodechar{≠}{\ensuremath{\neq}}
\newunicodechar{≤}{\ensuremath{\leq}}
\newunicodechar{≥}{\ensuremath{\geq}}
\newunicodechar{∀}{\ensuremath{\forall}}
\newunicodechar{∃}{\ensuremath{\exists}}
\newunicodechar{∘}{\ensuremath{\circ}}
\newunicodechar{·}{\ensuremath{\cdot}}
\newunicodechar{∈}{\ensuremath{\in}}
\newunicodechar{∉}{\ensuremath{\notin}}
\newunicodechar{≃}{\ensuremath{\simeq}}
\newunicodechar{≅}{\ensuremath{\cong}}
\newunicodechar{≈}{\ensuremath{\approx}}
\newunicodechar{∧}{\ensuremath{\wedge}}
\newunicodechar{∨}{\ensuremath{\vee}}
\newunicodechar{¬}{\ensuremath{\neg}}
\newunicodechar{⊢}{\ensuremath{\vdash}}
\newunicodechar{⟨}{\ensuremath{\langle}}
\newunicodechar{⟩}{\ensuremath{\rangle}}
\newunicodechar{⇑}{\ensuremath{\mathord{\uparrow}}}
\newunicodechar{↑}{\ensuremath{\mathord{\uparrow}}}
\newunicodechar{∎}{\ensuremath{\blacksquare}}
\newunicodechar{∣}{\ensuremath{\mid}}
\newunicodechar{∤}{\ensuremath{\nmid}}
\newunicodechar{∑}{\ensuremath{\textstyle\sum}}
\newunicodechar{∏}{\ensuremath{\textstyle\prod}}
\newunicodechar{∅}{\ensuremath{\emptyset}}
\newunicodechar{∪}{\ensuremath{\cup}}
\newunicodechar{∩}{\ensuremath{\cap}}
\newunicodechar{≡}{\ensuremath{\equiv}}
\newunicodechar{ℕ}{\ensuremath{\mathbb{N}}}
\newunicodechar{ℤ}{\ensuremath{\mathbb{Z}}}
\newunicodechar{ℝ}{\ensuremath{\mathbb{R}}}
\newunicodechar{𝔽}{\ensuremath{\mathbb{F}}}
\newunicodechar{α}{\ensuremath{\alpha}}
\newunicodechar{β}{\ensuremath{\beta}}
\newunicodechar{σ}{\ensuremath{\sigma}}
\newunicodechar{τ}{\ensuremath{\tau}}
\newunicodechar{φ}{\ensuremath{\varphi}}
\newunicodechar{ψ}{\ensuremath{\psi}}
\newunicodechar{χ}{\ensuremath{\chi}}
\newunicodechar{Φ}{\ensuremath{\Phi}}
\newunicodechar{Δ}{\ensuremath{\Delta}}
\newunicodechar{₀}{\textsubscript{0}}
\newunicodechar{₁}{\textsubscript{1}}
\newunicodechar{₂}{\textsubscript{2}}
\newunicodechar{ᵏ}{\ensuremath{{}^{k}}}
\newunicodechar{ⁿ}{\ensuremath{{}^{n}}}
\newunicodechar{ʲ}{\ensuremath{{}^{j}}}
\newunicodechar{²}{\ensuremath{{}^{2}}}
\newunicodechar{³}{\ensuremath{{}^{3}}}
\newunicodechar{⁴}{\ensuremath{{}^{4}}}
\newunicodechar{⁻}{\ensuremath{{}^{-}}}
\newunicodechar{¹}{\ensuremath{{}^{1}}}

% ---- Lean listing style ----
\definecolor{leanbg}{RGB}{247,247,250}
\definecolor{leankw}{RGB}{0,90,160}
\definecolor{leancm}{RGB}{120,120,120}
\lstdefinestyle{lean}{
  backgroundcolor=\color{leanbg},
  basicstyle=\ttfamily\small,
  keywordstyle=\color{leankw}\bfseries,
  commentstyle=\color{leancm}\itshape,
  morekeywords={theorem,lemma,def,fun,Prop,Type,variable,where,
    Field,Fintype,Finite,DecidableEq,CommRing,CommSemiring,Ring,CharP,Function,
    Injective,Surjective,Bijective,by,rfl,have,rw,exact,ext,simp,intro,
    iterateFrobenius,RingHom,Commute,convert,apply,using,Nat,Coprime,gcd,
    Odd,Even,IsAPN,IsAB,walsh,Nat,obtain,induction,cases,calc,grind},
  showstringspaces=false,
  columns=fullflexible,
  extendedchars=true,
  frame=single,
  framesep=6pt,
  xleftmargin=4pt,
  aboveskip=12pt,
  belowskip=14pt,
}
\lstset{style=lean}

\newtheorem{lemma}{Lemma}
\newtheorem{theorem}{Theorem}
\newtheorem{corollary}{Corollary}

% boxed "what just happened" note
\newcommand{\note}[1]{%
  \par\smallskip
  \noindent\colorbox{leanbg}{\parbox{\dimexpr\linewidth-2\fboxsep}{\small #1}}%
  \par\smallskip}


% --- Local overrides: A4, comfortable margins, open spacing ---
\geometry{a4paper,margin=2.6cm}
\onehalfspacing
\setlength{\emergencystretch}{8em}
\setlength{\parskip}{0.6em}

% Tighten the boxed-note macro for the streamlined layout
\renewcommand{\note}[1]{%
  \par\smallskip
  \noindent\colorbox{leanbg}{\parbox{\dimexpr\linewidth-2\fboxsep}{\small #1}}%
  \par\smallskip}

% Highlight macros: Big Ideas and Cool Facts
\definecolor{ideabg}{RGB}{233,243,255}
\definecolor{ideabar}{RGB}{0,90,160}
\definecolor{factbg}{RGB}{236,249,238}
\definecolor{factbar}{RGB}{30,140,70}

\newcommand{\bigidea}[1]{%
  \par\medskip
  \noindent\begin{tikzpicture}
    \node[fill=ideabg,rounded corners,inner sep=8pt,
          text width=\dimexpr\linewidth-2.2cm\relax] (b)
      {\small\textbf{\textcolor{ideabar}{Big idea.}}\;\; #1};
  \end{tikzpicture}\par\medskip}

\newcommand{\coolfact}[1]{%
  \par\medskip
  \noindent\begin{tikzpicture}
    \node[fill=factbg,rounded corners,inner sep=8pt,
          text width=\dimexpr\linewidth-2.2cm\relax] (b)
      {\small\textbf{\textcolor{factbar}{Cool fact.}}\;\; #1};
  \end{tikzpicture}\par\medskip}

\title{\textbf{Differential Uniformity, One Closure Law:\\[2pt]
the Frobenius Twist and the Kasami Even-$k$ Move}}
\author{}
\date{}

\begin{document}

\maketitle
\thispagestyle{empty}

\begin{abstract}
\noindent
\emph{Differential uniformity} $\delta(f)$ measures how well an S-box
resists differential cryptanalysis. The two best classes are APN
($\delta=2$) and PN ($\delta=1$).

This pearl rests on one law: \emph{relabelling the output of $f$ by an
additive bijection leaves $\delta$ unchanged.} From it we get two
payoffs, almost for free:
\begin{itemize}
  \item the \textbf{Frobenius twist} --- raising the output to a power of
        two keeps APN;
  \item the \textbf{Kasami even-$k$ move} --- the classical odd-$k$ APN
        theorem implies the even-$k$ case, with no new analysis.
\end{itemize}
Every boxed statement is checked in Lean~4 / Mathlib. One law about
addition does all the work.
\end{abstract}

\section{Introduction}

Differential cryptanalysis attacks a cipher by feeding it input pairs
that differ by a fixed amount $a$, then watching which output
difference $b$ shows up most often. An S-box that smears differences out
evenly is hard to attack.

The measure is the \emph{differential uniformity} $\delta(f)$: the most
input pairs that share any one nonzero input difference and output
difference. Smaller is better. Two classes sit at the bottom:
\[
  \text{APN: } \delta(f)=2 \quad(\text{char } 2),
  \qquad
  \text{PN: } \delta(f)=1 \quad(\text{odd char}).
\]
Power functions $x \mapsto x^d$ supply most known APN families: Gold,
Kasami, Welch, Niho, inverse, Dobbertin.

\bigidea{We do not re-prove a hard APN theorem. We isolate \emph{one}
elementary law --- $\delta$ ignores additive relabelling of the output
--- and show how much it buys. The hard classical input (odd-$k$ Kasami)
is a black box; the new content is the \emph{reduction}.}

The trick is to state the definition with \emph{no field structure},
using only $+$ and $-$. Then the closure law is one line, and the
Frobenius corollary is just a special case.

\section{The invariant}

Push $f$ sideways by a step $a$ and see how much it moved:
\[
  \text{derivative of } f \text{ in direction } a:
  \qquad f(x+a)-f(x).
\]

Fix a nonzero step $a$ and a target $b$, then count the hits:
\[
  \mathrm{fiberCard}\, f\, a\, b
  \;=\;
  \#\{\, x : f(x+a) - f(x) = b \,\}.
\]

Take the worst case over all $a\neq 0$ and all $b$:
\[
  \delta(f) \;=\; \max_{a \neq 0,\; b}\; \mathrm{fiberCard}\, f\, a\, b.
\]

Flat differences mean small $\delta$, which means hard to attack:
\[
  \boxed{\;\mathrm{IsAPN}\, f \;:\Longleftrightarrow\; \delta(f) = 2\;}
  \qquad
  \boxed{\;\mathrm{IsPN}\, f \;:\Longleftrightarrow\; \delta(f) = 1\;}
\]

\note{The definition uses only $+$ and $-$, never multiplication. So we
state it for \emph{any} finite additive groups $V,W$ and $f:V\to W$ --- no
field, no characteristic. That generality is what makes the closure law
free.}

\subsection{The primitive, in Lean}

\begin{lstlisting}
def derivMap (f : V → W) (a x : V) : W := f (x + a) - f x

def fiberCard (f : V → W) (a : V) (b : W) : ℕ :=
  (univ.filter fun x => derivMap f a x = b).card

def differentialUniformity (f : V → W) : ℕ :=
  (univ.filter fun a => a ≠ 0).sup
    fun a => univ.sup fun b => fiberCard f a b

def IsAPN (f : V → W) : Prop := differentialUniformity f = 2
\end{lstlisting}

\noindent
Four lines, and \texttt{V}, \texttt{W} are any finite additive groups.
Using \texttt{Finset.sup} keeps $\delta$ a natural number, so the
closure-law rewrites stay purely combinatorial.

\coolfact{In characteristic two, ``$\delta(f)\le 2$'' has a hands-on
shape: among any three inputs sharing a derivative value, two must
coincide. ``At most two-to-one.''}

\begin{theorem}[collision form]
If $\mathrm{derivMap}\, f\, a\, x = \mathrm{derivMap}\, f\, a\, y
= \mathrm{derivMap}\, f\, a\, z$, then $x=y$, $x=z$, or $y=z$.
\end{theorem}

\noindent
The two descriptions are proved equal in Lean
(\texttt{isAtMostTwoToOne\_iff\_diffUnif\_le\_two}).

\subsection{A worked example}

Take $F=\mathrm{GF}(8)$ and the Gold map $f(x)=x^{3}$ (the smallest
Kasami exponent, $d_1=2^2-2^1+1$). Fix $a\neq 0$. In characteristic two,
\[
  f(x+a)-f(x) = (x+a)^3 - x^3 = a\,(x^2 + a x) + a^3 .
\]
For a target $b$, the equation $f(x+a)-f(x)=b$ is a quadratic in $x$.
Over characteristic two such a quadratic has at most two roots, and
$x\mapsto x+a$ leaves it fixed --- so roots come in pairs $\{x,x+a\}$.

Every nonempty fibre therefore has exactly two elements:
$\mathrm{fiberCard}\,f\,a\,b \in\{0,2\}$, so $\delta(f)=2$. This is the
base case the Kasami reduction rests on.

\section{The one closure law}

This is the heart of the paper. Let $\sigma$ be an \emph{additive
bijection} of the output (an element of \texttt{W ≃+ W'}): a relabelling
that respects $+$.

\begin{theorem}[closure law]
$\delta(\sigma \circ f) = \delta(f).$
\end{theorem}

\noindent
Relabelling the output cannot change how spread the differences are.
Here is the whole reason. First, $\sigma$ slides through the derivative:
\[
  \mathrm{derivMap}\,(\sigma \circ f)\, a\, x
  = \sigma\big(\mathrm{derivMap}\, f\, a\, x\big),
\]
because $\sigma(u-v)=\sigma(u)-\sigma(v)$. Second, it just renames each
fibre one-for-one:
\[
  \mathrm{fiberCard}\,(\sigma \circ f)\, a\, b
  = \mathrm{fiberCard}\, f\, a\, (\sigma^{-1} b).
\]
Same counts, shuffled. Same maximum.

\begin{lstlisting}
lemma derivMap_comp_addEquiv (f : V → W) (σ : W ≃+ W') (a x : V) :
    derivMap (σ ∘ f) a x = σ (derivMap f a x)

lemma fiberCard_comp_addEquiv (f : V → W) (σ : W ≃+ W') (a) (b) :
    fiberCard (σ ∘ f) a b = fiberCard f a (σ.symm b)

lemma differentialUniformity_comp_addEquiv
    (f : V → W) (σ : W ≃+ W') :
    differentialUniformity (σ ∘ f) = differentialUniformity f

theorem IsAPN.comp_addEquiv
    {f : V → W} (hf : IsAPN f) (σ : W ≃+ W') : IsAPN (σ ∘ f)
\end{lstlisting}

\bigidea{One move carries a whole equivalence class. Prove a single
example APN, and get every additive relabelling for free. The law never
sees that $W$ might be a field.}

\section{Payoff one: the Frobenius twist}

Now work in a finite field $F$ of characteristic two. The Frobenius
power $\sigma : x \mapsto x^{2^{j}}$ is
\begin{itemize}
  \item \textbf{additive:} $(x+y)^{2^j} = x^{2^j} + y^{2^j}$ (the
        freshman's dream, iterated);
  \item \textbf{bijective:} an injective endomorphism of a finite field.
\end{itemize}
So it is just one instance of the closure law.

\begin{corollary}
If $f$ is APN, then $x \mapsto f(x)^{2^{j}}$ is APN.
\end{corollary}

\begin{lstlisting}
theorem IsAPN.comp_frobenius
    {f : F → F} (hf : IsAPN f) (j : ℕ) :
    IsAPN (fun x => (f x) ^ (2 ^ j))
\end{lstlisting}

\noindent
The proof: build $\sigma$ from Mathlib's \texttt{iterateFrobenius}, then
apply \texttt{IsAPN.comp\_addEquiv}. Raising the output to a power of two
is invisible to APN-ness --- and that freedom drives the second payoff.

\section{Payoff two: the Kasami even-$k$ move}

The Kasami power function on $\mathrm{GF}(2^n)$ is $x \mapsto x^{d_k}$,
\[
  d_k = 2^{2k} - 2^{k} + 1 .
\]
Classically it is APN whenever $\gcd(k,n)=1$, proved via the \emph{odd}
$k$ case. We handle even $k$ not by new analysis, but by \emph{moving}
the odd case onto it.

First, restate the law as a plain bijection, and specialise $\sigma$ to a
Frobenius power on $f=x^d$:

\begin{lstlisting}
lemma apn_comp_additive_bijective
    {f : F → F} (hf : IsAPN f)
    {σ : F → F} (hσ_bij : Function.Bijective σ)
    (hσ_add : ∀ x y, σ (x + y) = σ x + σ y) : IsAPN (σ ∘ f)

lemma apn_frob_twist
    {d : ℕ} (hd : IsAPN (fun x : F => x ^ d)) (j : ℕ) :
    IsAPN (fun x : F => x ^ (d * 2 ^ j))
\end{lstlisting}

\noindent
Multiplying the \emph{exponent} by $2^j$ keeps APN, since
$x^{d \cdot 2^j} = (x^{d})^{2^j}$ is a Frobenius twist.

\subsection{The key congruence}

On $\mathrm{GF}(2^n)$ the unit group has order $2^n-1$, so for $x\neq 0$
exponents only matter modulo $2^n-1$ (and $x=0$ is fixed by every power
map). The Kasami exponents at $k$ and at its complement $n-k$ are linked:

\begin{lemma}[exponent congruence]
For $0 < k < n$, \;$d_k \equiv d_{\,n-k}\cdot 2^{\,2k} \pmod{2^{n}-1}$.
\end{lemma}

\noindent
A short computation:
\[
  d_{n-k}\cdot 2^{2k}
  = (2^{2(n-k)}-2^{n-k}+1)\,2^{2k}
  = 2^{2n}-2^{n+k}+2^{2k} .
\]
Reduce $2^{2n}\equiv 2^n\equiv 1$ and $2^{n+k}\equiv 2^k$ modulo
$2^n-1$, and what remains is $2^{2k}-2^k+1 = d_k$. So on the field,
\[
  x^{\,d_k} = \big(x^{\,d_{\,n-k}}\big)^{2^{\,2k}} :
\]
one map, written either as a Kasami exponent or as a Frobenius twist of
another (\texttt{kasami\_pow\_frob\_identity}).

\begin{lemma}[APN of the complement]
If $x^{\,d_{\,n-k}}$ is APN, then $x^{\,d_k}$ is APN.
\end{lemma}

\begin{lstlisting}
lemma kasami_apn_of_complement
    (hn : Fintype.card F = 2 ^ n)
    (k : ℕ) (hk : 0 < k) (hkn : k < n)
    (hapn : IsAPN (fun x : F => x ^ kasamiExp (n - k))) :
    IsAPN (fun x : F => x ^ kasamiExp k)
\end{lstlisting}

\subsection{Parity flips for free}

Why is the complement the right partner? Let $n$ be \emph{odd} and $k$ be
\emph{even} with $\gcd(k,n)=1$. Then:
\begin{itemize}
  \item $n-k = \text{odd}-\text{even}$ is \emph{odd};
  \item coprimality survives: $\gcd(n-k,n)=\gcd(k,n)=1$.
\end{itemize}
So $x^{\,d_{\,n-k}}$ is covered by the classical odd theorem, and the
complement trick lifts it back to $x^{\,d_k}$.

\coolfact{The even case is the odd case seen through Frobenius. The
parity barrier simply dissolves --- no second analytic proof needed.}

\begin{lstlisting}
theorem kasami_is_apn_even_k
    (hn : Fintype.card F = 2 ^ n) (k : ℕ)
    (hk : 1 < k) (hk_even : Even k) (hkn : k < n)
    (hn_odd : Odd n) (hcop : Nat.Coprime k n)
    (hnk : n - k ≥ 2) :
    IsAPN (fun x : F => x ^ kasamiExp k)

theorem kasami_is_apn_general
    (hn : Fintype.card F = 2 ^ n) (k : ℕ)
    (hk : 1 < k) (hkn : k < n)
    (hn_odd : Odd n) (hcop : Nat.Coprime k n) :
    IsAPN (fun x : F => x ^ kasamiExp k)
\end{lstlisting}

\noindent
The edge $k=n-1$ lands on the Gold exponent $d_1=3$, proved APN directly.
With it, \texttt{kasami\_is\_apn\_general} covers every $k$ coprime to
$n$, with no parity restriction left.

\section{The bridge back}

We started field-free; we end where the textbooks are.

\begin{theorem}
In characteristic two, if $\delta(f) \le 2$ then for $a \neq 0$,
\[
  f(x+a) - f(x) = f(y+a) - f(y)
  \;\Longrightarrow\;
  y = x \;\text{ or }\; y = x + a.
\]
\end{theorem}

\noindent
Nothing is lost: the classical collision statement is \emph{derived},
not assumed (\texttt{atMostTwoToOne\_charTwo\_collision}). The group
definition and the field-theoretic folklore coincide exactly.

\section{Formalization notes}

A few choices keep things short and reusable:
\begin{itemize}
  \item \textbf{$\delta$ over groups, not fields.} The primitive uses
        only $+$ and $-$, so the closure law is a fact about
        \texttt{AddEquiv} and Frobenius is a special case. The field
        enters only to make Frobenius additive and bijective.
  \item \textbf{$\delta$ valued in \texttt{ℕ}.} Using \texttt{Finset.sup}
        of \texttt{fiberCard} makes APN-ness a decidable equation and
        fibre re-indexing a \texttt{Finset} bijection.
  \item \textbf{Exponents mod $2^n-1$.} The Kasami identity is proved as
        an equality of functions on $F$ (split on $x=0$, use the order of
        $F^\times$), avoiding congruence side-conditions.
  \item \textbf{Odd case as a hypothesis.} The even-$k$ theorem takes the
        odd-$k$ result as input, so it is robust to whichever proof of
        the odd case is supplied.
\end{itemize}
The files are \texttt{DifferentialUniformity.lean} (definition, collision
form, closure law, Frobenius twist) and \texttt{KasamiEvenK.lean}
(congruence, complement trick, even-$k$ and general theorems).

\section{Related work}

APN functions come from differential cryptanalysis and have a large
literature of power-function families: Kasami $2^{2k}-2^k+1$ with
$\gcd(k,n)=1$, Gold $2^k+1$, and others. The odd-$k$ APN property is
textbook; the even-$k$ reduction through the complement is folklore. Our
contribution is to make that reduction fully explicit and
machine-checked, with the closure law as its only structural ingredient.
Mathlib supplies the finite-field and Frobenius infrastructure
(\texttt{frobenius}, \texttt{iterateFrobenius}, order of the unit group);
what is new is the small, self-contained theory of $\delta$, stated for
arbitrary finite additive groups.

\section{Conclusion}

The take-away is methodological. One almost trivial law --- $\delta$ is
blind to additive relabelling of the output --- is enough to
\begin{enumerate}
  \item record the Frobenius twist for free, and
  \item collapse the parity split in Kasami to one exponent congruence
        plus the same law.
\end{enumerate}
Stating the invariant at the right level of generality, with no field
structure, is what makes both payoffs mechanical.

\bigidea{Big idea, small footprint: one law about addition does all the
work.}

\end{document}
